\begin{table*}[!h]
    \centering
    \small
    \caption{RQ3 (Why): Topic-level mechanisms for seller deception/collusion (LLM-labeled).}
    \label{tab:rq3_collusion_mechanisms}
    \begin{tabular}{p{0.12\textwidth}p{0.25\textwidth}p{0.55\textwidth}}
    \toprule
    \textbf{Topic} & \textbf{Label} & \textbf{Mechanism Explanation} \\
    \midrule
    T0 & Mixed Quality Strategy & Seller deception emerges as sellers strategically mix high-quality (HQ) and low-quality (LQ) products to maximize profits while managing their reputation. By advertising HQ products with truth warrants and LQ products without, they aim to appeal to different buyer segments while minimizing the risk of negative feedback. \\
    T1 & Profit Maximization & Seller deception emerges as individuals prioritize profit over ethical considerations, strategically listing low-quality products as high-quality to enhance profit margins while attempting to maintain a positive reputation through the use of Truth Warrants. This creates a conflict between the desire for financial gain and the obligation to provide truthful representations to buyers. \\
    \bottomrule
    \end{tabular}
\end{table*}